\documentclass[a4paper,11pt]{book}
%\documentclass[a4paper,twoside,11pt,titlepage]{book}
\usepackage{listings}
\usepackage[utf8]{inputenc}
\usepackage[spanish]{babel}

% \usepackage[style=list, number=none]{glossary} %
%\usepackage{titlesec}
%\usepackage{pailatino}

\decimalpoint
\usepackage{dcolumn}
\newcolumntype{.}{D{.}{\esperiod}{-1}}
\makeatletter
\addto\shorthandsspanish{\let\esperiod\es@period@code}
\makeatother


%\usepackage[chapter]{algorithm}
\RequirePackage{verbatim}
%\RequirePackage[Glenn]{fncychap}
\usepackage{fancyhdr}
\usepackage{graphicx}
\usepackage{afterpage}

\usepackage{longtable}

\usepackage[pdfborder={000}]{hyperref} %referencia

% ********************************************************************
% Re-usable information
% ********************************************************************
\newcommand{\myTitle}{Análisis de los algoritmos de Boosting en entornos con ruido de clase.\xspace}
\newcommand{\myDegree}{Doble Grado en Ingeniería Informática y Matemáticas\xspace}
\newcommand{\myName}{Rafael Nogales Vaquero (alumno)\xspace}
\newcommand{\myProf}{Francisco Herrera Triguero  (tutor1)\xspace}
\newcommand{\myOtherProf}{Salvador García López (tutor2)\xspace}
%\newcommand{\mySupervisor}{Put name here\xspace}
\newcommand{\myFaculty}{Escuela Técnica Superior de Ingenierías Informática y de
Telecomunicación\xspace}
\newcommand{\myFacultyShort}{E.T.S. de Ingenierías Informática y de
Telecomunicación\xspace}
\newcommand{\myDepartment}{Departamento de ...\xspace}
\newcommand{\myUni}{\protect{Universidad de Granada}\xspace}
\newcommand{\myLocation}{Granada\xspace}
\newcommand{\myTime}{\today\xspace}
\newcommand{\myVersion}{Version 0.1\xspace}


\hypersetup{
pdfauthor = {\myName (email (en) ugr (punto) es)},
pdftitle = {\myTitle},
pdfsubject = {},
pdfkeywords = {palabra_clave1, palabra_clave2, palabra_clave3, ...},
pdfcreator = {LaTeX con el paquete ....},
pdfproducer = {pdflatex}
}

%\hyphenation{}


%\usepackage{doxygen/doxygen}
%\usepackage{pdfpages}
\usepackage{url}
\usepackage{colortbl,longtable}
\usepackage[stable]{footmisc}
%\usepackage{index}

%\makeindex
%\usepackage[style=long, cols=2,border=plain,toc=true,number=none]{glossary}
% \makeglossary

% Definición de comandos que me son tiles:
%\renewcommand{\indexname}{Índice alfabético}
%\renewcommand{\glossaryname}{Glosario}

\pagestyle{fancy}
\fancyhf{}
\fancyhead[LO]{\leftmark}
\fancyhead[RE]{\rightmark}
\fancyhead[RO,LE]{\textbf{\thepage}}
\renewcommand{\chaptermark}[1]{\markboth{\textbf{#1}}{}}
\renewcommand{\sectionmark}[1]{\markright{\textbf{\thesection. #1}}}

\setlength{\headheight}{1.5\headheight}

\newcommand{\HRule}{\rule{\linewidth}{0.5mm}}
%Definimos los tipos teorema, ejemplo y definición podremos usar estos tipos
%simplemente poniendo \begin{teorema} \end{teorema} ...
\newtheorem{teorema}{Teorema}[chapter]
\newtheorem{ejemplo}{Ejemplo}[chapter]
\newtheorem{definicion}{Definición}[chapter]

\definecolor{gray97}{gray}{.97}
\definecolor{gray75}{gray}{.75}
\definecolor{gray45}{gray}{.45}
\definecolor{gray30}{gray}{.94}

\lstset{ frame=Ltb,
     framerule=0.5pt,
     aboveskip=0.5cm,
     framextopmargin=3pt,
     framexbottommargin=3pt,
     framexleftmargin=0.1cm,
     framesep=0pt,
     rulesep=.4pt,
     backgroundcolor=\color{gray97},
     rulesepcolor=\color{black},
     %
     stringstyle=\ttfamily,
     showstringspaces = false,
     basicstyle=\scriptsize\ttfamily,
     commentstyle=\color{gray45},
     keywordstyle=\bfseries,
     %
     numbers=left,
     numbersep=6pt,
     numberstyle=\tiny,
     numberfirstline = false,
     breaklines=true,
   }
 
% minimizar fragmentado de listados
\lstnewenvironment{listing}[1][]
   {\lstset{#1}\pagebreak[0]}{\pagebreak[0]}

\lstdefinestyle{CodigoC}
   {
	basicstyle=\scriptsize,
	frame=single,
	language=C,
	numbers=left
   }
\lstdefinestyle{CodigoC++}
   {
	basicstyle=\small,
	frame=single,
	backgroundcolor=\color{gray30},
	language=C++,
	numbers=left
   }

 
\lstdefinestyle{Consola}
   {basicstyle=\scriptsize\bf\ttfamily,
    backgroundcolor=\color{gray30},
    frame=single,
    numbers=none
   }


\newcommand{\bigrule}{\titlerule[0.5mm]}


%Para conseguir que en las páginas en blanco no ponga cabecerass
\makeatletter
\def\clearpage{%
  \ifvmode
    \ifnum \@dbltopnum =\m@ne
      \ifdim \pagetotal <\topskip
        \hbox{}
      \fi
    \fi
  \fi
  \newpage
  \thispagestyle{empty}
  \write\m@ne{}
  \vbox{}
  \penalty -\@Mi
}
\makeatother

\usepackage{pdfpages}
\begin{document}
\begin{titlepage}
 
 
\newlength{\centeroffset}
\setlength{\centeroffset}{-0.5\oddsidemargin}
\addtolength{\centeroffset}{0.5\evensidemargin}
\thispagestyle{empty}

\noindent\hspace*{\centeroffset}\begin{minipage}{\textwidth}

\centering
\includegraphics[width=0.9\textwidth]{imagenes/logo_ugr.jpg}\\[1.4cm]

\textsc{ \Large TRABAJO FIN DE GRADO\\[0.2cm]}
\textsc{ DOBLE GRADO EN INGENIERÍA INFORMÁTICA Y MATEMÁTICAS }\\[1cm]
% Upper part of the page
% 
% Title
{\Huge\bfseries Análisis de los algoritmos de Boosting en entornos con ruido de clase.\\
}
\noindent\rule[-1ex]{\textwidth}{3pt}\\[3.5ex]
{\large\bfseries Rafael Nogales Vaquero}
\end{minipage}

\vspace{2.5cm}
\noindent\hspace*{\centeroffset}\begin{minipage}{\textwidth}
\centering

\textbf{Autor}\\ {Rafael Nogales Vaquero}\\[2.5ex]
\textbf{Directores}\\
{Francisco Herrera Triguero \\
Julián Luengo Martín}\\[2cm]
\includegraphics[width=0.3\textwidth]{imagenes/etsiit_logo.png}\\[0.1cm]
\textsc{Escuela Técnica Superior de Ingenierías Informática y de Telecomunicación}\\
\textsc{---}\\
Granada, 7 Septiembre de 2018
\end{minipage}
%\addtolength{\textwidth}{\centeroffset}
%\vspace{\stretch{2}}
\end{titlepage}



\chapter*{}
%\thispagestyle{empty}
%\cleardoublepage

%\thispagestyle{empty}

\begin{titlepage}
 
 
\setlength{\centeroffset}{-0.5\oddsidemargin}
\addtolength{\centeroffset}{0.5\evensidemargin}
\thispagestyle{empty}

\noindent\hspace*{\centeroffset}\begin{minipage}{\textwidth}

\centering
%\includegraphics[width=0.9\textwidth]{imagenes/logo_ugr.jpg}\\[1.4cm]

%\textsc{ \Large PROYECTO FIN DE CARRERA\\[0.2cm]}
%\textsc{ INGENIERÍA EN INFORMÁTICA}\\[1cm]
% Upper part of the page
% 

 \vspace{3.3cm}

%si el proyecto tiene logo poner aquí
%\includegraphics{imagenes/logo.png} 
% \vspace{0.5cm}

% Title

{\Huge\bfseries Análisis de los algoritmos de Boosting en entornos con ruido de clase.\\
}
\noindent\rule[-1ex]{\textwidth}{3pt}\\[3.5ex]
%{\large\bfseries .\\[4cm]}
\end{minipage}

\vspace{2.5cm}
\noindent\hspace*{\centeroffset}\begin{minipage}{\textwidth}
\centering

\textbf{Autor}\\ {Rafael Nogales Vaquero}\\[2.5ex]
\textbf{Directores}\\
{Francisco Herrera Triguero\\
Julian Luengo Martín}\\[2cm]
%\includegraphics[width=0.15\textwidth]{imagenes/tstc.png}\\[0.1cm]
%\textsc{Departamento de Teoría de la Señal, Telemática y Comunicaciones}\\
%\textsc{---}\\
%Granada, mes de 201
\end{minipage}
%\addtolength{\textwidth}{\centeroffset}
\vspace{\stretch{2}}

 
\end{titlepage}






\cleardoublepage
\thispagestyle{empty}

\begin{center}
{\large\bfseries Análisis de los algoritmos de Boosting en entornos con ruido de clase.}\\
\end{center}
\begin{center}
Rafael Nogales Vaquero \\
\end{center}

%\vspace{0.7cm}
\noindent{\textbf{Palabras clave}: Machine Learning, Boosting, Ruido de clase, Función de pérdida, Problemas de Clasificación}\\

\vspace{0.7cm}
\noindent{\textbf{Resumen}}\\

En este TFG se propone cambiar la función de pérdida de los algoritmos de tipo boosting con el objetivo de hacerlos más robustos a conjuntos de datos de entrenamiento parcialmente mal clasificados.\cleardoublepage


\thispagestyle{empty}


\begin{center}
{\large\bfseries Analysis of Boosting algorithms in environments with class noise.}\\
\end{center}
\begin{center}
Rafael Nogales Vaquero\\
\end{center}

%\vspace{0.7cm}
\noindent{\textbf{Keywords}: Machine Learning, Boosting, Class noise, Loss function, Classification problems}\\

\vspace{0.7cm}
\noindent{\textbf{Abstract}}\\

In this TFG it is proposed to change the loss function of the boosting algorithms in order to make them more robust to partially misclassified training data sets.

\chapter*{}
\thispagestyle{empty}

\noindent\rule[-1ex]{\textwidth}{2pt}\\[4.5ex]

Yo, \textbf{Rafael Nogales Vaquero}, alumno de la titulación DOBLE GRADO EN INGENIERÍA INFORMÁTICA Y MATEMÁTICAS de la \textbf{Escuela Técnica Superior
de Ingenierías Informática y de Telecomunicación de la Universidad de Granada}, con DNI 76669550Q, autorizo la
ubicación de la siguiente copia de mi Trabajo Fin de Grado en la biblioteca del centro para que pueda ser
consultada por las personas que lo deseen.

\vspace{6cm}

\noindent Fdo: Rafael Nogales Vaquero

\vspace{2cm}

\begin{flushright}
Granada a 7 de Agosto de 2018
\end{flushright}


\chapter*{}
\thispagestyle{empty}

\noindent\rule[-1ex]{\textwidth}{2pt}\\[4.5ex]

D. \textbf{Francisco Herrera Triguero  }, Profesor del Departamento de Ciencias de la Computación e Inteligencia Artificial  de la Universidad de Granada.

\vspace{0.5cm}

D. \textbf{Julian Luengo Martín}, Profesor del Departamento de Ciencias de la Computación e Inteligencia Artificial  de la Universidad de Granada.


\vspace{0.5cm}

\textbf{Informan:}

\vspace{0.5cm}

Que el presente trabajo, titulado \textit{\textbf{Título del proyecto, Subtítulo del proyecto}},
ha sido realizado bajo su supervisión por \textbf{Rafael Nogales Vaquero }, y autorizamos la defensa de dicho trabajo ante el tribunal
que corresponda.

\vspace{0.5cm}

Y para que conste, expiden y firman el presente informe en Granada a X de mes de 201 .

\vspace{1cm}

\textbf{Los directores:}

\vspace{5cm}

\noindent \textbf{Francisco Herrera Triguero   \ \ \ \ \ Julian Luengo Martín }

\chapter*{Agradecimientos}
\thispagestyle{empty}

       \vspace{1cm}


A mi familia, amigos y profesores


%\frontmatter
%\tableofcontents
%\listoffigures
%\listoftables
%
%\mainmatter
%\setlength{\parskip}{5pt}

%\chapter{Introducción}

\section{Contexto}

A día de hoy hay un gran interés en la inteligencia artificial en general y particularmente existe una gran demanda de modelos predictivos de todo tipo en la mayoría de empresas de mediano y gran tamaño, incluso en algunas startups de base tecnológica.  

En este trabajo vamos a interesarnos por las diferentes formas en las que podemos evitar una de las mayores dificultades a las que las organizaciones se encuentran a la hora de de construir sus modelos predictivos: los datos de mala calidad (en particular los datos con ruido).

Por otro lado en plataformas de Data Science se ha visto que el algoritmo XGBoost ha sido el más elegido para ganar numerosas competiciones de construcción de modelos predictivos. Por ello vamos a analizar cómo utilizar éste algoritmo de la forma más robusta posible, es decir evitando que sus predicciones se vean afectadas por datos de entrenamiento de baja calidad.

Este algoritmo es utilizado en problemas de ámbitos muy diferentes siempre que haya que trabajar con datos tabulados suele dar un gran resultado.  
Por desgracia hay muchos problemas con datos tabulados que son susceptibles de tener datos mal clasificados y esto es un gran problema para los algoritmos de boosting clásicos como Adaboost.  

El problema principal radica en que para ajustarse rápidamente a los datos los algoritmos de boosting suelen utilizar funciones de pérdida exponenciales o cuadráticas que dan una importancia enorme a los datos ruidosos. 

Lo cual nos deja un amplio terreno sobre el que trabajar.

\section{Solución}
Debido a que nuestro mayor problema es el sobreajuste causado al aprender de los datos ruidosos vamos a utilizar técnicas que den la menor importancia posible a aquellos datos que considere ruidosos.  
Es por esto que la estrategia a seguir para mitigar este problema va a consistir en utilizar otras funciones de pérdida más robustas al ruido en combinación con técnicas de regularización.  
Previamente repasaremos todos los conceptos matemáticos relacionados con el boosting, las funciones de pérdida y el aprendizaje automático en general. 

\section{Herramientas}
Clasificaremos nuestras herramientas en dos tipos: Herramientas matemáticas y herramientas software.  

Para la parte matemática nuestras herramientas serán el cálculo diferencial, la estadística y el álgebra lineal.

En la parte informática utilizaremos principalmente R, Python y algunas de sus librerías: pandas, scikit-learn, xgboost, matplotlib...


%
%\input{capitulos/02_EspecificacionRequisitos}
%
%\input{capitulos/03_Planificacion}
%
%\input{capitulos/04_Analisis}
%
%\input{capitulos/05_Diseno}
%
%\input{capitulos/06_Implementacion}
%
%\input{capitulos/07_Pruebas}
%
%\input{capitulos/08_Conclusiones}
%
%%\chapter{Conclusiones y Trabajos Futuros}
%
%
%%\nocite{*}
%\bibliography{bibliografia/bibliografia}\addcontentsline{toc}{chapter}{Bibliografía}
%\bibliographystyle{miunsrturl}
%
%\appendix
%\input{apendices/manual_usuario/manual_usuario}
%%\input{apendices/paper/paper}
%\input{glosario/entradas_glosario}
% \addcontentsline{toc}{chapter}{Glosario}
% \printglossary
\chapter*{}
\thispagestyle{empty}

\end{document}
